\section{Conclusion}
\label{sec:conclusion}

This paper centered on a core question: while preserving the Tucker low-rank backbone as much as possible, how can one introduce effective nonlinearity in a controlled manner? To answer this question, we proposed NTD-PL. The method takes a shared Tucker latent tensor as the structural backbone and imposes an explicit polynomial link at the entry level, thereby compressing nonlinear enhancement into only a few link coefficients. NTD-PL places greater emphasis on structural continuity, controllable parameterization, and interpretability of the underlying mechanism.

Around this modeling idea, we provided a complete argument from model to experiments. Theoretically, NTD-PL has a clear nesting relation with standard Tucker: the linear special case degenerates exactly to Tucker, while increasing the degree yields a controlled model expansion. It achieves exact matching under polynomial generation and finite-order approximation with truncation error bounds under analytic generation. Algorithmically, we presented a unified learning framework that can handle both full-observation decomposition and masked completion, and that remains consistent with practical implementation.

The experimental results support these conclusions. In the linear-consistency experiments, restricted NTD-PL almost coincides with Tucker, showing that the model does not fabricate artificial gains in purely linear regimes. In controlled nonlinear residual experiments, the performance trends are consistent with the theoretical expectations. Real hyperspectral experiments further show that, on both full-observation reconstruction and random-missing completion tasks on CAVE, NTD-PL consistently outperforms Tucker under matched parameter budgets, and the advantage is consistent at both the scene-wise and rate-wise levels.

More importantly, the mechanism and spatial evidence explain why the method is effective. The mechanism comparison shows that post-hoc polynomial calibration on Tucker outputs can only partially approximate NTD-PL and cannot replace the joint modeling of nonlinear links with the low-rank backbone. The spatial analysis shows that the gain is not globally uniform smoothing, but is mainly concentrated in difficult local regions with strong boundaries. Degree analysis further shows that the main gains are concentrated in low to medium orders and that marginal improvements gradually saturate as the degree increases, indicating that the additional computational cost of the method is quantifiable and bounded. Cross-dataset results also lead to a clear conclusion: NTD-PL is effective on most real scenes, but not unconditionally so. When the data are closer to linear mixing and the entry-level nonlinear deviation is weaker, its additional gains are correspondingly reduced.

Future work may proceed along three directions. First, one may investigate more stable or more flexible explicit link families, such as orthogonal bases or piecewise bases, to improve numerical conditioning at higher orders. Second, one may develop more efficient mask-aware training and approximate computation to improve large-scale completion efficiency. Third, one may carry out more systematic theoretical analysis of identifiability, generalization error, and adaptive degree selection. Overall, NTD-PL provides an analyzable, implementable, and extensible path for ``enhancing Tucker,'' and offers a reusable starting point for structured nonlinear low-rank modeling.
